
\section{Method}
\label{sec:method}
%\fixme{in general, the methodology is kinda short, and do not make it read like an engineering work, try to add more scientific/algorithmic flavor. The best way is to introduce variables/equations....In the current section, i did not even see any equation......}
The overall framework of the proposed AnyTalker is illustrated in~\cref{fig:method}. 
AnyTalker inherits certain architectural components from the Wan I2V model~\cite{wan2025wan}.
To handle multi-stream audio and identity input, we introduce a specialized multi-stream processing structure, termed as the \textit{Audio-Face Cross Attention}~(AFCA), which will be further described in~\cref{subsec:audio-face-attn}. 
Our training pipeline is divided into two stages, which are summarized in~\cref{subsec:training-strategy}.

\begin{figure}[t]
    \centering
    \includegraphics[width=0.9\linewidth]{fig/attn_mask_face_mask.pdf}
    \caption{(a) Mapping of video tokens to audio tokens, facilitated by a custom attention mask. 
    Every 4 audio tokens are bound to 1 video token, except for the first. 
    (b) Mask token used for output masking in Audio-Face Cross Attention.}
    \label{fig:token}
\end{figure}

\subsection{Preliminaries}
As a DiT-based model, AnyTalker tokenizes the 3D VAE features $f_{video}$ through patchifying and flattening, whereas the text features $f_{text}$ are generated by the T5 encoder~\cite{raffel2020exploring}. 
Additionally, AnyTalker incorporates Reference Attention Layer, a cross-attention mechanism that leverages the CLIP image encoder~\cite{radford2021learning} \(E_{\text{CLIP}} \) to extract features $f_{ref}$ from the first frame of the video. Wav2Vec2~\cite{baevski2020wav2vec} is also applied to extract the audio feature $f_{audio}$. 
The overall input features $f_{input}$ can be written as
\begin{equation}
    f_{input} = [f_{video}, f_{text}, f_{ref}, f_{audio}].
\end{equation}
\label{eq:feature}
Consistent with the Wan model, all attention layers are connected to the final output FFN layer (omitted in~\cref{fig:method}).

\subsection{Audio-Face Cross Attention}
\label{subsec:audio-face-attn}
To enable multi-person talking, the model must be capable of handling multi-stream audio inputs.
Potential solutions may include the L-RoPE technique used in MultiTalk~\cite{kong2025let}, which assigns unique labels and biases to different audio features.
However, the range of these labels needs to be explicitly defined, which limits its scalability.
Considering this, we design a more extensible structure to drive multiple IDs and enable accurate control in a scalable manner.
As depicted in~\cref{fig:method}~(a) and (c), we introduce a specialized structure named Audio-Face Cross Attention (AFCA). 
The structure can iterate through a loop multiple times, contingent upon the number of input face-audio pairs. 
As illustrated in~\cref{eq:attn} and~\cref{fig:method} (c), it enables flexible processing of diverse audio and identity inputs, with summed outputs from each iteration yielding the final attention output.

\noindent \textbf{Audio Token Modeling.}
We employ Wav2Vec2~\cite{baevski2020wav2vec} to encode audio features. 
The first latent frame attends to all audio tokens, whereas each subsequent latent frame focuses only on a local temporal window corresponding to four audio tokens. 
This structured alignment between video and audio streams is achieved by applying a Temporal Attention Mask \( M_{\text{temporal}} \), as explicitly shown in~\cref{fig:token}~(a). 
Furthermore, to enable comprehensive information integration, each audio token \( f_{\text{audio}} \) utilized in the AFCA computation is concatenated with a face token \( f_{\text{face}} \) encoded by \( E_{\text{CLIP}} \).
This concatenation allows all video query tokens \( Q_{\text{video}} \) to attend to different pairs of audio and face information effectively, as computed below:
\begin{equation}
\label{eq:af-attention}
\begin{aligned}
    K_{\text{af}} &= \text{Concat}(f_{\text{audio}}, f_{\text{face}})\cdot W_K, \\
    V_{\text{af}} &= \text{Concat}(f_{\text{audio}}, f_{\text{face}})\cdot W_V, \\
    \text{Attn}_{\text{out}} &= \text{MHCA}(Q_{\text{video}}, K_{\text{af}}, V_{\text{af}}, M_{\text{temporal}}).
\end{aligned}
\end{equation}
Here, MHCA denotes \textit{Multi-Head Cross Attention}, while $W_K$ and $W_V$ represent the key matrix and value matrix, respectively.
The attention output $\text{Attn}_{\text{out}}$ will later be refined by the face mask token, as described in~\cref{eq:mask}.


\noindent \textbf{Face Token Modeling.}
The facial image is obtained by online cropping the first frame of the selected video clip using InsightFace~\cite{deng2020retinaface} during training, while the facial mask $M_{\text{face}}$ is precomputed offline to cover the maximum extent of the face mask across the entire video, i.e., the \textit{global} face bounding box. 
This mask ensures that facial movements will never exceed this region, preventing the mask from incorrectly activating video tokens after the reshape and flatten operations shown in~\cref{fig:token}~(b), especially for videos with significant facial displacements.
This mask, which shares the same dimensions as \( \text{Attn}_{\text{out}} \), can be directly employed for element-wise multiplication to compute the \textit{Audio-Face Cross Attention} output, as formulated below:
\begin{equation}
\begin{aligned}
{M}_\text{token} &= \text{Patchify}( \text{Flatten}(M_\text{face})), \\
\text{AFCA}_\text{out} &= {M}_\text{token} \odot \text{Attn}_{\text{out}}.  \label{eq:mask}
\end{aligned}
\end{equation}
% Consequently, the forward pass of the AnyTalker Attention Block, preceding the FFN, is summarized as:
Consequently, the hidden state $H_{i}$ of each I2V DiT block can be formulated as
\begin{equation}
\begin{aligned}
H_{i}^{'} &= H_{i} + \text{AFCA}_\text{out}^{(1)}
+ \cdots + \text{AFCA}_\text{out}^{(n)}, \label{eq:attn}
\end{aligned} 
\end{equation}
where $i$ represents the layer index of the attention block, and $n$ denotes the number of IDs.
Note that all $\mathrm{AFCA}_{\mathrm{out}}$ terms are produced by the \emph{same} AFCA layer with shared parameters. 
The AFCA computation is applied $n$ times iteratively, once for each individual.
This architecture enables the number of drivable IDs to scale arbitrarily.

\begin{figure*}[t]
    \centering
    \includegraphics[width=\linewidth]{fig/interactivity.pdf}
    \caption{
    Two video clips from InteractiveEyes with ${Motion}$ score (px): left shows original video, right shows cropped face and eye landmarks. Head turn toward the speaker or eyebrow raise will increase $Motion$ and Interactivity; sustained stillness keeps both low.
    }
    \label{fig:interactivity}
\end{figure*}

\subsection{Training Strategy}
\label{subsec:training-strategy}
AnyTalker explores the potential of single-person data for learning multi-person speaking patterns, with low-cost single-person data comprising the majority of training data.

\noindent \textbf{Single-Person Data Pretraining.}
We train the model using both standard single-person data and synthetic two-person data generated by horizontal concatenation. 
With an equal 50\% probability, each batch of data is randomly configured to either two-person or single-person mode, as depicted in~\cref{fig:method}~(b).
In two-person mode, each sample within the batch is horizontally concatenated with the data at the next index, along with its corresponding audio. 
This approach keeps the batch size identical across the two modes for every data batch.
Additionally, we have predefined several generic text prompts that describe dual people speaking when data concatenation happens.

Although the above data construction strategy enhances the model's ability to localize audio-visual features within local regions and learn speaking patterns of dual speakers, completely omitting the single-person data is not feasible. 
Doing so would significantly degrade the model's performance on generating accurate lip movements, leading to unstable driving results, which we later discuss in~\cref{tab:ab-data}.

% \noindent \textbf{Multi-person Data Refinement.}
% In the second stage, we fine-tune the model using pure multi-person data to enhance the interactivity within different IDs in the generated videos.
% Although our training data contains only interactions between two identities, we find that the model naturally generalizes to scenarios with more than two IDs, as shown in~\cref{fig:teaser}.
% We speculate that this is because the AFCA mechanism has truly learned the general patterns of human interaction, including not only the accurate lip-syncing to the audio but also the listening and responsive behaviors to other IDs' speaking actions.

% For the creation of the multi-person data, we conducted rigorous quality control, using InsightFace~\cite{deng2019arcface} to ensure two faces in most frames, audio diarization~\cite{Plaquet23} to separate audio for two or one speakers, optical flow~\cite{karaev2024cotracker} to filter excessive motion, and Sync scores~\cite{chung2016out} to pair audio with faces. 
% This pipeline yielded 12 hours of high-quality two-person data, significantly less than previous methods~\cite{kong2025let, wang2025interacthuman, ma2025playmate2}. 
% More details are in the supplementary material.
% As the design of AnyTalker's AFCA Layer inherently supports multi-ID inputs, two-person data is fed into the model in the same format as the concatenated data in the first stage, with no extra processing required.
%
%yc
\noindent \textbf{Multi-person Data Refinement.}
In the next stage, we refine the model using a small amount of authentic multi-person data to enhance the interactivity within different IDs.
Although our training data contains only interactions between two identities, we surprisingly find that our model equipped with the AFCA module naturally generalizes to scenarios with more than two IDs, as shown in~\cref{fig:teaser}.
We speculate that this is because the AFCA mechanism enables learning general patterns of human interaction, including not only accurate lip-syncing to the audio but also listening and responsive behaviors to other IDs' speaking actions.

To construct high-quality multi-person training data, we construct a rigorous quality control pipeline, using InsightFace~\cite{deng2019arcface} to ensure two faces in most frames, audio diarization~\cite{Plaquet23} to separate audio and ensure there is only one or two speakers, optical flow~\cite{karaev2024cotracker} to filter excessive motion, and Sync scores~\cite{chung2016out} to pair audio with faces. 
More details about this pipeline are in the supplementary material.
This pipeline yields a total of $12$ hours of high-quality dual-person data, which is a small amount compared to previous methods~\cite{kong2025let, wang2025interacthuman, ma2025playmate2}. 
As the design of AnyTalker's AFCA Layer inherently supports multi-ID inputs, two-person data is fed into the model in the same format as the concatenated data in the first stage, with no extra processing required.

%yc
To summarize, the single-person data training process enhances the model's lip-syncing capability and generation quality, while also learning a generalized multi-person speaking pattern. Subsequently, lightweight multi-person data refinement compensates for the real interactions that cannot be fully covered by single-person data.



